%!TEX root = ./../../thesis-trier.tex

The preceding sections have linked the absence of certain congruences to synchronizability, reachability, and distinguishability
properties. Here, we collect some sufficient conditions for the absence of non-trivial $S$-congruences
in a transformation semigroups~$S$.

For $g = (0 ~ 1 ~ \cdots n - 1)$ and $t \colon [n] \to [n]$
let~$r$ be the greatest common divisor of the numbers $t(i) - i$
for $i \in [n]$ and~$n$.
If $r > 1$, then the mod-$r$ residue classes form a non-trivial $\langle G \cup \{t\} \rangle_S$-congruence.
Next, we show that under certain assumptions,
$r = 1$ also implies that there exist no non-trivial congruences.
This statement extends a statement from~\cite[Lem.~3]{HeathIKS09}.

\begin{proposition} \label{syncmax:primitive:prop:circular-primitive}
	Let $g = (0 ~ 1 ~ \cdots ~ n - 1)$
	and $t_1, \ldots, t_k \colon [n] \to [n]$
	be transformations moving at least one point such that,
	for each $i \in \{1,\ldots,k\}$, we have
	$n \ge 2 M_i - 1$
	or $|\move{t_i}| < d$ for all proper divisors~$d$ of~$n$ with $d > 1$
	where $M_i= \max(\move{t_i})$.
	If the greatest common divisor
	of the numbers $t_i(j) - j$
	for $j \in \move{t_i}$ and $i \in \{1,\ldots,k\}$
	is coprime to~$n$,
	then $S = \langle \{g,t_1,\ldots,t_k\} \rangle_S$
	admits no non-trivial $S$-congruence.
\end{proposition}
\begin{proof}
	Assume there exists an $S$-congruence
	that does not equal the universal relation.
	Since this is also a $\langle g \rangle_S$-congruence,
	by \cref{lem:congruence_on_cyclic_group}, each class
	is a mod-$r$ residue class where~$r$ divides~$n$
	and we have $r < n$. For a contradiction, assume $r > 1$.
	
	Next, we show that for each $i \in \{1,\ldots,k\}$, 
	each congruence class contains an element fixed by~$t_i$. 
	To see this, note that for $|\move{t_i}| \ge n/r$
	we have $n - M_i \ge r$ since $r \le n/2$ as it is a proper divisor
	and, with the assumption $n \ge 2 M_i - 1$, we find
	\[
	n - M_i
	\ge n - \frac{n + 1}{2}
	= \frac{n - 1}{2}
	\ge r - \frac{1}{2}.
	\]
	Thus, if $|\move{t_i}| \ge n / r$,
	then since $\{M_i, M_i + 1, \ldots, n\} \subseteq \fix{t_i}$
	and $n - M_i \ge r$ each congruence class
	contains a point from $\fix{t_i}$, and
	if $|\move{t_i}| < n / r$, since each congruence class has size $n / r$,
	it also contains a point from $\fix{t_i}$.
	
	Next, consider $j \in \move{t_i}$.
	As written above, there exists $l \in \fix{t_i}$
	with $j \equiv l \pmod{r}$.
	Thus we have $t_i(l) \equiv l \pmod{r}$
	and since $\equiv$ is an $S$-congruence
	and $j \equiv l \pmod{r}$
	we get $t_i(j) \equiv j \pmod{r}$, i.e.,
	$r$ divides $t_i(j) - j$.
	Hence, $r$ divides the numbers $t_i(j) - j$ for $i \in \{1,\ldots,k\}$
	and $j \in \move{t_i}$.
	But~$r$ also divides $n$, and by the assumption,
	it follows that $r = 1$, giving a contradiction.	
\end{proof}

\begin{corollary}
	Let $g = (0 ~ 1 ~ \cdots ~ n - 1)$
	and $t_1, \ldots, t_k \colon [n] \to [n]$
	be simple idempotents. If the greatest common divisor
	of the numbers $t_i(j) - j$ 
	for $j \in [n]$ with $t_i(j) \ne j$
	where $i \in \{1,\ldots,k\}$
	is coprime to~$n$, then $S = \langle \{g,t_1,\ldots,t_k\} \rangle_S$
	admits no non-trivial $S$-congruence.	
\end{corollary}
\begin{proof}
	This follows from \cref{syncmax:primitive:prop:circular-primitive}
	since for simple idempotents we have $|\move{t_i}| = 1$.
\end{proof}

\begin{example}
	Let $G = \langle \{ (0 ~ 1 ~ \cdots ~ n - 1) \} \rangle_G$.
	Then for non-empty $A \subseteq [n]$ where there exists
	no divisor $r$ of~$d$ with $r > 1$ such that $i \equiv j \pmod{r}$
	for all $i,j \in A$
	we have that $G$ is $A$-primitive
	by \cref{lem:congruence_on_cyclic_group}.
	In this case, by \cref{syncmax:primitive:prop:circular-primitive},
	when adding simple idempotents mapping
	every two distinct states in~$A$ onto each other,
	the generated transformation semigroup~$S$
	does not admit any non-trivial $S$-congruences.
	For example, if~$n$ is even and $n \ge 8$,
	then $G$ is $\{0,2,5,7\}$-primitive.
	Adding a simple idempotent mapping $0$ to~$5$
	results in a transformation semigroup admitting
	no non-trivial congruence,
	but adding the two simple idempotents
	mapping $0$ to~$2$ 
	and $5$ to~$7$
	results in a transformation semigroup
	that admits a congruence whose classes
	are the mod-$2$ residue classes.
	By \cref{syncmax:primitive:prop:circular-primitive} we have to add simple idempotents $t$
	such that the numbers $t(i) - i$
	and~$n$ are coprime. In fact, for simple idempotents,
	this condition characterizes those sets
	of simple idempotents
	resulting in transformation semigroups
	admitting a non-trivial congruence.	
\end{example}

\begin{lemma} \label{syncmax:primitive:lem:transitive-primitive}
	Let~$G$ be a transitive permutation group on~$X$ and~$t$ be a simple idempotent
	with $t(a) \ne a$ and $S = \langle G \cup \{t\} \rangle_S$.
	Then $t(a) \sim a$ for every non-trivial $S$-congruence $\sim$
	on~$X$.	
\end{lemma}
\begin{proof}
	Let $\Delta$ be an equivalence class for some non-trivial $S$-congruence.
	Since this is also a $G$-congruence, $\Delta$ is a block of~$G$
	of size at least two.
	By transitivity, we have $a \in \Delta^g$ for some $g \in G$.
	Let $b \in \Delta^g$ with $b \ne a$.
	We have $b \sim a$.
	Since~$\sim$ is a $S$-congruence, we find $b = t(b) \sim t(a)$
	and so $t(a) \sim a$.	
\end{proof}

\begin{corollary}
	Let~$G$ be a transitive and $A$-primitive permutation group on~$X$
	for some $A \subseteq X$.
	For a set $T \subseteq \mathcal T_X$
	of simple idempotents
	such that $(X, E)$ with $E = \{ \{a, t(a)\} \mid a \in X, t(a) \ne a \}$
	is a connected graph,
	the semigroup $S = \langle G \cup T \rangle_S$
	admits no non-trivial $S$-congruences.
\end{corollary}
\begin{proof}
	By \cref{syncmax:primitive:lem:transitive-primitive}
	and the assumption, the points in~$A$
	are all congruent for each $S$-congruence
	not equal to the equality relation,
	and thus by $A$-primitivity
	the universal relation is the only
	$S$-congruence not equal to the equality relation.	
\end{proof}

\begin{lemma}
	Let $G$ be a permutation group on~$X$
	having a non-empty set of fixpoints~$F$
	and where $X \setminus F$ is an orbit.
	Then the only blocks containing~$F$
	are $F$ and~$X$.		
\end{lemma}
\begin{proof}
	Clearly,~$F$ is a block.
	Next, let $\Delta$
	be a block properly containing~$F$.
	Pick $a \in \Delta \setminus F$.
	Assume there exists $b \in X \setminus \Delta$.
	By assumption, we have $g(a) = b$
	for some $g \in G$.
	Now, $F \subseteq g(\Delta)$.
	Thus, since~$\Delta$
	is a block, we have $\Delta = g(\Delta)$.
	However, since $b \in g(\Delta)$, this gives $b \in \Delta$,
	which is a contradiction.
\end{proof}

Let $\sim$ be an $S$-congruence.
Clearly, $t( [a]_{\sim} ) \subseteq [a]_{\sim}$ for $a \in \fix{t}$.
Thus, if for a transformation semigroup $S$ on~$X$ with $S$-congruence $\sim$
and $A \subseteq X$, for each pair of distinct $a,b \in X$
there exist $t \in S$ and $c \in \fix{t}$ with $t(a) = b$ and $c \sim a$,
we have %$c = t(c) \sim t(a) = b$
$b = t(a) \subseteq [a]_{\sim}$
and this yields that all points in~$A$ are congruent w.r.t.~$\sim$.

\begin{proposition}
	Let $G$ be $\{a,b\}$-primitive for distinct $a, b$ on a set $X$ of size at least three
	and $S = \langle G \cup \{t\} \rangle_S$
	for a simple idempotent $t \colon X \to X$ mapping $a$ to $b$.
	Then~$G$ is transitive and $S$ admits
	no non-trivial $S$-congruence,
	or~$G$ is intransitive 
	with two distinct orbits $\Orbit{a}{G}$
	and $\Orbit{b}{G}$
	and the non-trivial $S$-congruences
	are precisely those induced by sets of the form $\{c\}$
	for $c \in \Orbit{a}{G}$
	and the equivalence classes 
	of some $H$-congruence on  $\Orbit{b}{G}$ for $H = \{ g_{|\Orbit{b}{G}} \mid g \in G \}$
	that does not equal the equality relation.
\end{proposition}
\begin{proof}
	Let $\sim$ be some $S$-congruence not equal to the equality relation.
	First, assume there exists $c \in X \setminus \{a\}$
	with $c \sim a$. Then $c = t(c) \sim t(a) = b$
	and so $a \sim b$, which yields that $\sim$
	is the universal relation.
	Otherwise, $\{a\}$ is a congruence class.
	If~$G$ is transitive, then all blocks have size one
	and so~$\sim$ equals the equality relation.
	If~$G$ is intransitive, 
	then by \cref{sync_max:prop:ab_non_congruent_then_two_orbits} 
	we have two disjoint orbits $\Orbit{a}{G}$
	and $\Orbit{b}{G}$.
	Since~$\{a\}$ is an equivalence class, and so a block for~$G$,
	we can conclude that all equivalence classes
	in $\Orbit{a}{G}$ are singletons.
	Thus, for $c \in \Orbit{b}{G}$
	we have $[c]_{\sim} \subseteq \Orbit{b}{G}$
	and since $\Orbit{b}{G} \subseteq \fix{t}$,
	$[c]_{\sim}$ is a block.
	Next, it is easy to verify that the given sets induce $S$-congruences on~$X$.
	So, we obtain the claim.		
\end{proof}

\begin{example}
	The group $\langle (0)(1 ~ 2 ~ 3) \rangle_G$
	is $\{0,1\}$-primitive, intransitive
	and the partition $\{0\}$, $\{1,2,3\}$
	induces a congruence for the transformation semigroup
	generated when adding the simple idempotent mapping $0$
	to~$1$.	
\end{example}